% Actividad Semana 2: Shiny frente a Power BI (una columna, A4, citas APA 7)
% Compilar:  tectonic -X compile informe.tex
\documentclass[12pt,a4paper]{article}
\usepackage[spanish,es-noshorthands,es-tabla,es-nodecimaldot]{babel}
\usepackage{fontspec}
\setmainfont{texgyretermes}[Extension=.otf,UprightFont=*-regular,BoldFont=*-bold,ItalicFont=*-italic,BoldItalicFont=*-bolditalic]
\setsansfont{texgyreheros}[Extension=.otf,UprightFont=*-regular,BoldFont=*-bold,ItalicFont=*-italic,BoldItalicFont=*-bolditalic]
\setmonofont{texgyrecursor}[Extension=.otf,UprightFont=*-regular,BoldFont=*-bold,ItalicFont=*-italic,BoldItalicFont=*-bolditalic,Scale=0.95]
\usepackage[a4paper,margin=2.5cm]{geometry}
\usepackage{setspace}
\usepackage{graphicx}
\usepackage{booktabs,tabularx,array,multirow}
\usepackage[table]{xcolor}
\usepackage{caption,subcaption}
\usepackage{float}
\usepackage{enumitem}
\usepackage[round,authoryear]{natbib}
\usepackage{titlesec}
\usepackage{fancyhdr}
\usepackage[hidelinks]{hyperref}
\usepackage{xurl}
\usepackage[normalem]{ulem}

\definecolor{marino}{HTML}{14213D}
\definecolor{azul}{HTML}{2F6FED}
\definecolor{suave}{HTML}{64748B}
\definecolor{fondo}{HTML}{EEF2F8}
\hypersetup{colorlinks=true,linkcolor=black,citecolor=black,urlcolor=black}

\onehalfspacing
\setlength{\emergencystretch}{3em}
\setlength{\parindent}{1.27cm}
\setlength{\parskip}{2pt}
\captionsetup{font=small,labelfont=bf,justification=justified,singlelinecheck=true}
\captionsetup[figure]{name=Figura}
\captionsetup[sub]{font=small,labelfont=bf}
\titleformat{\section}{\large\bfseries}{\thesection.}{0.6em}{}
\titleformat{\subsection}{\normalsize\bfseries}{\thesubsection}{0.6em}{}
\titlespacing*{\section}{0pt}{14pt}{6pt}
\titlespacing*{\subsection}{0pt}{10pt}{4pt}
\setlist{itemsep=2pt,topsep=4pt}

\setlength{\headheight}{15pt}
\pagestyle{fancy}
\fancyhf{}
\fancyhead[L]{\small Actividad Semana 2: Herramientas de analítica y \emph{reporting}}
\fancyhead[R]{\small Shiny y Power BI}
\fancyfoot[C]{\small\thepage}
\renewcommand{\headrulewidth}{0.4pt}

% Página de capturas: Shiny arriba y Power BI abajo
\newcommand{\pareja}[6]{% 1 sh, 2 pbi, 3 alto shiny, 4 pie shiny, 5 pie pbi, 6 título+etiqueta
  \begin{figure}[p]\centering
  \begin{subfigure}{\textwidth}\centering
    \includegraphics[height=#3,width=\textwidth,keepaspectratio]{figs2/#1}
    \caption{#4}
  \end{subfigure}\par\medskip
  \begin{subfigure}{\textwidth}\centering
    \includegraphics[width=0.97\textwidth]{figs2/#2}
    \caption{#5}
  \end{subfigure}
  #6
  \end{figure}}

\raggedbottom
\makeatletter
\setlength{\@fptop}{0pt}
\setlength{\@fpsep}{12pt}
\setlength{\@fpbot}{0pt plus 1fil}
\makeatother

% Enlaces con la letra del texto, en azul y subrayados
\newcommand{\enlace}[2]{\href{#1}{\textcolor{blue}{#2}}}

\begin{document}

% =============================================================================
\thispagestyle{plain}
\begin{center}
{\large\bfseries Shiny y Power BI en la analítica básica y el \emph{reporting} de datos:\\
órdenes de compra pública de la región Puno\par}
\vspace{10pt}
{Mary Luz Nina Palacios\par}
{\small Universidad Nacional del Altiplano, Escuela de Posgrado. Maestría, curso Business Intelligence and
Data Analytics. Docente: Fred Torres Cruz\par}
{\small Actividad Semana 2: Herramientas de analítica y \emph{reporting}. Puno, octubre de 2026\par}
\end{center}
\vspace{4pt}

% =============================================================================
\section*{Resumen}
\noindent El presente informe compara Shiny y Power BI a partir de un caso real: las 5\,925 órdenes de compra que
las entidades públicas de la región Puno emitieron en los Catálogos Electrónicos de Perú Compras entre
enero de 2025 y septiembre de 2026, por un total de S/~85,19 millones. Con el mismo archivo de datos se
construyó un módulo de análisis exploratorio en cada herramienta, con seis secciones equivalentes, y se
comprobó que ambos producen exactamente las mismas cifras. Posteriormente, las herramientas se calificaron en doce
criterios con una escala de 1 a 5 cuyos niveles están definidos de antemano, de modo que cada puntaje
responde a algo observable durante el desarrollo. Shiny obtuvo 53 de 60 puntos y Power BI 47. La diferencia
proviene de la analítica básica (24 frente a 18), donde R aporta histogramas, densidades y mapas con
límites distritales; en el \emph{reporting} ambas sumaron 29 puntos, aunque por razones distintas. Power BI
ofrece interacción sin necesidad de programar y una menor curva de aprendizaje, mientras que Shiny permite
un diseño libre y una publicación abierta sin costo. Los resultados no señalan una herramienta superior,
sino dos herramientas complementarias cuya elección depende del perfil de los usuarios del reporte.

\medskip\noindent\textbf{Palabras clave:} inteligencia de negocios, análisis exploratorio de datos,
tableros, Shiny, Power BI, compras públicas.


% =============================================================================
\section{Introducción}
Disponer de datos no garantiza, por sí solo, mejores decisiones. Su valor depende de que se conviertan en
información comprensible y oportuna para quien debe decidir, y esa es precisamente la tarea de la inteligencia de negocios
\citep{watson2007,chen2012}. En la práctica, este trabajo comprende dos momentos. Primero se exploran los datos para
entender cómo se comportan, detectar asimetrías o valores atípicos y formular buenas preguntas
\citep{behrens1997}. Luego se comunican los resultados en reportes o tableros que otras personas consultan
de manera autónoma \citep{yigitbasioglu2012,sarikaya2019}. La herramienta que se use en cada momento influye en
lo que se puede analizar, en su costo y en la facilidad para mantener el trabajo.

Shiny y Power BI representan dos maneras distintas de resolver ese problema. Shiny es un paquete de R con
el que se construyen aplicaciones web interactivas escribiendo código, y se usa cada vez más para poner
análisis estadísticos al alcance de usuarios que no programan \citep{jia2022}. Power BI es la plataforma de
inteligencia de negocios de Microsoft, pensada para preparar datos, modelarlos y publicar reportes desde una
interfaz visual \citep{becker2019}. Ambas permiten explorar datos y construir tableros, pero parten de
lógicas distintas: la primera se basa en el código y la segunda en una interfaz de arrastrar y soltar.

Para que la comparación tuviera validez práctica se descartó un conjunto de datos de ejemplo y se trabajó con
información real, reciente y de la región: las compras que municipalidades, el Gobierno Regional,
hospitales, universidades y otras entidades de Puno hicieron a través de los Catálogos Electrónicos de
Perú Compras. Estos datos abiertos tienen valor propio, porque permiten conocer en qué se gasta el
presupuesto público, a qué proveedores se adquiere y cómo cambió el gasto de un año a otro \citep{janssen2012}. De hecho, los registros de órdenes y
contratos públicos ya se han usado para construir indicadores que alertan sobre riesgos en las compras del
Estado \citep{fazekas2020}.
Así, cada herramienta se puso a prueba con preguntas de negocio y no con ejercicios de ejemplo.

El objetivo del trabajo es comparar ambas herramientas como medios de analítica básica y de
\emph{reporting}, desarrollando en cada una un módulo de análisis exploratorio sobre los mismos datos y
calificándolas con criterios explícitos. En las secciones siguientes se describen los datos, se presentan
los dos módulos sección por sección, se resumen los hallazgos, se explica cómo se asignó cada puntaje y,
al final, se discuten los resultados.

% =============================================================================
\section{Datos y preparación}
Los datos provienen del conjunto \emph{Órdenes de compra realizadas a través de los Catálogos
Electrónicos}, que Perú Compras publica mes a mes en la Plataforma Nacional de Datos Abiertos
\citep{perucompras2026}. Se descargaron los 21 archivos mensuales de enero de 2025 a septiembre de 2026.
Cada archivo trae las órdenes de todo el país con la entidad compradora, el proveedor, el rubro o acuerdo
marco, el tipo de procedimiento, el estado de la orden y el monto.

La fuente no indica la región de la entidad, así que se obtuvo del padrón reducido del RUC de la SUNAT
\citep{sunat2026}: a cada entidad se le asignó el ubigeo de su domicilio fiscal y se conservaron solo las
del departamento de Puno (código 21). El mismo padrón permitió saber dónde tiene su domicilio cada empresa
proveedora y, con ello, qué parte del gasto se destina a empresas de la región. Para el mapa se usaron los
límites distritales del INEI. Toda esta preparación se hizo una sola vez, en Python, y ambas
herramientas partieron del mismo archivo, lo que asegura que cualquier diferencia entre ellas venga de la
herramienta y no de los datos.

El archivo final tiene 5\,929 órdenes y 25 variables. Antes de analizarlo se revisó su calidad
(Tabla~\ref{tab:calidad}). No hubo órdenes duplicadas, valores faltantes en las variables clave ni montos
iguales o menores a cero. Cuatro órdenes estaban resueltas o anuladas, es decir, no llegaron a
ejecutarse, y se excluyen por defecto; por eso el análisis trabaja con 5\,925 órdenes.

\begin{table}[H]\centering\small
\caption{Revisión de la calidad de los datos}\label{tab:calidad}
\begin{tabular}{lr}
\toprule
\textbf{Verificación} & \textbf{Resultado} \\ \midrule
Órdenes en el archivo & 5\,929 \\
Órdenes duplicadas & 0 \\
Valores faltantes en variables clave & 0 \\
Montos iguales o menores a cero & 0 \\
Órdenes resueltas o anuladas (excluidas por defecto) & 4 \\
Proveedores sin domicilio fiscal publicado en la SUNAT & 126 \\ \bottomrule
\end{tabular}
\end{table}

% =============================================================================
\section{Los dos módulos de reporte}
\subsection{Módulo en Shiny}
La aplicación está en un solo archivo, \texttt{app.R}, de unas 1\,390 líneas. Usa \texttt{bslib} para el
diseño, \texttt{plotly} para los gráficos, \texttt{leaflet} y \texttt{sf} para el mapa
\citep{pebesma2018} y \texttt{DT} para las tablas. A la izquierda hay un panel de filtros (periodo,
sector, provincia, rubro y estado) y en la parte superior, seis pestañas: Resumen, Proveedores, Entidades, Mapa
distrital, Exploración y Datos. Los filtros se aplican con un botón, para que la aplicación no recalcule
todo cada vez que se marca una opción, y otro botón los devuelve a su estado inicial. Cada tarjeta de
indicador muestra la variación frente a 2025 y, al pasar el cursor, explica en lenguaje sencillo qué
significa la cifra. La aplicación está publicada en \enlace{https://mlnp.shinyapps.io/compras-puno/}{\uline{https:/}\allowbreak \uline{/}\allowbreak \uline{mlnp.}\allowbreak \uline{shinyapps.}\allowbreak \uline{io/}\allowbreak \uline{compras-}\allowbreak \uline{puno/}} y su
código en \enlace{https://github.com/mlnp-dev/compras-puno-shiny}{\uline{https:/}\allowbreak \uline{/}\allowbreak \uline{github.}\allowbreak \uline{com/}\allowbreak \uline{mlnp-}\allowbreak \uline{dev/}\allowbreak \uline{compras-}\allowbreak \uline{puno-}\allowbreak \uline{shiny}}.

\subsection{Módulo en Power BI}
En Power BI se cargó el mismo archivo y se armó un modelo con dos tablas relacionadas, una de órdenes y
otra de distritos con sus coordenadas, además de una tabla pequeña para elegir la métrica de los gráficos.
Los cálculos están en 54 medidas DAX: indicadores, variaciones frente al año anterior, concentración,
Pareto, estadísticos descriptivos y textos de ayuda. El reporte tiene las mismas seis secciones, con un
menú lateral, filtros desplegables en cada página, un selector que cambia todos los gráficos entre monto,
órdenes o empresas, una ventana emergente con la tendencia mensual y un botón para limpiar los filtros. Se
guardó en formato de proyecto (PBIP), que conserva el modelo y las páginas como archivos de texto. Ambos
módulos llevan arriba el mismo membrete institucional. El reporte se publicó en el servicio de Power BI
con una cuenta institucional y la opción de publicación en la web, de modo que puede consultarse sin
iniciar sesión en
\enlace{https://app.powerbi.com/view?r=eyJrIjoiMmJlMGJlMjQtZmU5My00YzU3LWEwZGYtMjdlNGJiZDMzNDMwIiwidCI6ImM0YTY2YzM0LTJiYjctNDUxZi04YmUxLWIyYzI2YTQzMDE1OCIsImMiOjR9}{\uline{https:/}\allowbreak \uline{/}\allowbreak \uline{app.}\allowbreak \uline{powerbi.}\allowbreak \uline{com/}\allowbreak \uline{view?r=eyJrIjoiM}\allowbreak \uline{mJlMGJlMjQtZmU5M}\allowbreak \uline{y00YzU3LWEwZGYtM}\allowbreak \uline{jdlNGJiZDMzNDMwI}\allowbreak \uline{iwidCI6ImM0YTY2Y}\allowbreak \uline{zM0LTJiYjctNDUxZ}\allowbreak \uline{i04YmUxLWIyYzI2Y}\allowbreak \uline{TQzMDE1OCIsImMiO}\allowbreak \uline{jR9}};
el proyecto completo está en \enlace{https://github.com/mlnp-dev/compras-puno-powerbi}{\uline{https:/}\allowbreak \uline{/}\allowbreak \uline{github.}\allowbreak \uline{com/}\allowbreak \uline{mlnp-}\allowbreak \uline{dev/}\allowbreak \uline{compras-}\allowbreak \uline{puno-}\allowbreak \uline{powerbi}}.

% =============================================================================
\section{Comparación sección por sección}
Las Figuras~\ref{fig:resumen} a \ref{fig:datos} muestran, para cada sección, la pantalla de Shiny arriba y
la de Power BI abajo, con los mismos datos y sin filtros. Verlas juntas permite notar dónde coinciden y
dónde cada herramienta resolvió la misma pregunta de otra manera.

\paragraph{Resumen.} Ambas pantallas presentan los mismos cuatro indicadores con su variación frente a
2025 (Figura~\ref{fig:resumen}). Shiny añade una comparación acumulada de 2026 contra 2025 y anotaciones
sobre el propio gráfico, como el pico de agosto o la caída de enero, que coincide con el inicio del año
fiscal. Power BI resuelve lo mismo con menos elementos, pero con su selector de métrica basta un clic para
que todas las figuras pasen de soles a número de órdenes o de empresas.

\paragraph{Proveedores.} Aquí se mide qué tan concentradas están las ventas (Figura~\ref{fig:proveedores}).
En Shiny la curva de Pareto se obtuvo ordenando y acumulando en dos líneas de R. En Power BI el mismo cálculo
exigió una medida con una tabla virtual, aunque el resultado se lee de forma muy directa en la tarjeta
\emph{152 de 712 empresas hacen el 80~\% del monto}.

\paragraph{Entidades.} Ambas herramientas muestran quién compra y en qué sector
(Figura~\ref{fig:entidades}). En Power BI, al hacer clic en una barra se filtra toda la página sin necesidad de
programación; en Shiny ese comportamiento tendría que programarse.

\paragraph{Mapa.} Es la sección donde más se notan las diferencias (Figura~\ref{fig:mapa}). Shiny dibuja los
límites reales de los 109 distritos con un color según el monto y, al hacer clic en uno, muestra quién
compra y en qué. Power BI ubica una burbuja por distrito sobre Azure Maps: la presentación es cuidada y el mapa
base es de buena calidad, pero no muestra los límites distritales y el visual depende de que el administrador de la
organización lo tenga habilitado. Además, el mapa de Power BI tardó varios segundos en dibujarse cada vez
que se abría la página. La solución de Shiny es la misma que se ha usado en tableros de vigilancia de salud
pública que combinan mapas y filtros \citep{lu2026}.

\paragraph{Exploración.} Es la sección más propia del análisis exploratorio (Figura~\ref{fig:exploracion}).
Shiny ofrece un histograma del monto en escala logarítmica, curvas de densidad por sector, un mapa de calor
por día de la semana y la tabla de valores atípicos. Power BI no tiene histograma ni gráfico de densidad
nativos, así que los montos se agruparon en cinco rangos y la media, la mediana, la desviación estándar, los
cuartiles y el coeficiente de variación se calcularon con medidas DAX. El resultado es correcto y fácil de
leer, pero obtenerlo requirió una solución alternativa.

\paragraph{Datos.} Ambas herramientas permiten revisar las órdenes una por una (Figura~\ref{fig:datos}).
Shiny trae un buscador y un botón que descarga en CSV solo lo filtrado; Power BI permite exportar desde el
menú del propio visual.

\pareja{sh_01_resumen.jpg}{pbi_01_resumen.jpg}{11.6cm}{Shiny, pestaña Resumen.}{Power BI, página Resumen.}
{\caption{Sección Resumen: indicadores con su variación frente a 2025, evolución mensual del gasto, rubros
y sectores.}\label{fig:resumen}}

\pareja{sh_02_empresas.jpg}{pbi_02_proveedores.jpg}{11.6cm}{Shiny, pestaña Proveedores.}{Power BI, página
Proveedores.}{\caption{Sección Proveedores: concentración de las ventas, origen de las empresas y relación
entre número de órdenes y monto.}\label{fig:proveedores}}

\pareja{sh_03_entidades.jpg}{pbi_03_entidades.jpg}{11.6cm}{Shiny, pestaña Entidades.}{Power BI, página
Entidades.}{\caption{Sección Entidades: entidades que más compran, sectores y provincias.}\label{fig:entidades}}

\pareja{sh_04_mapa.jpg}{pbi_04_mapa.jpg}{11.6cm}{Shiny, mapa con los límites de los distritos
(\texttt{leaflet} y \texttt{sf}).}{Power BI, burbujas por distrito sobre Azure Maps.}{\caption{Sección
Mapa: en ambas herramientas el distrito de Puno concentra el 71~\% del monto.}\label{fig:mapa}}

\pareja{sh_05_exploracion.jpg}{pbi_05_exploracion.jpg}{11.6cm}{Shiny, pestaña Exploración.}{Power BI,
página Exploración.}{\caption{Sección Exploración: estadísticos descriptivos, distribución de los montos,
días de emisión y órdenes atípicas.}\label{fig:exploracion}}

\pareja{sh_06_datos.jpg}{pbi_06_datos.jpg}{9cm}{Shiny, pestaña Datos.}{Power BI, página Datos.}
{\caption{Sección Datos: detalle de las órdenes con filtros y exportación.}\label{fig:datos}}

% =============================================================================
\section{Lo que muestran los datos}
Antes de comparar las herramientas es necesario verificar que ambas produzcan los mismos resultados. Las medidas de Power BI se
consultaron directamente en su motor de cálculo y se cotejaron con los valores de la aplicación Shiny.
Coincidieron en todos los indicadores (Tabla~\ref{tab:cifras}), de modo que las diferencias que se
discuten después son de forma y de proceso, no de resultado.

\begin{table}[H]\centering\small
\caption{Indicadores calculados en ambas herramientas}\label{tab:cifras}
\begin{tabular}{lrr}
\toprule
\textbf{Indicador} & \textbf{Shiny} & \textbf{Power BI} \\ \midrule
Monto total & S/ 85,19 millones & S/ 85,19 millones \\
Órdenes & 5\,925 & 5\,925 \\
Ticket promedio & S/ 14\,377 & S/ 14\,377 \\
Mediana del monto por orden & S/ 3\,370 & S/ 3\,370 \\
Empresas proveedoras & 712 & 712 \\
Variación del monto, enero a septiembre de 2026 frente a 2025 & $-$22,1~\% & $-$22,1~\% \\
Variación del número de órdenes en el mismo periodo & $-$11,8~\% & $-$11,8~\% \\
Monto que va a empresas de Puno & 12,5~\% & 12,5~\% \\
Empresas que concentran el 80~\% del monto & 152 & 152 \\ \bottomrule
\end{tabular}
\end{table}

El primer resultado relevante es la reducción de las compras en 2026. Entre enero y septiembre el monto fue
22,1~\% menor que en los mismos meses de 2025 (S/~29,60 millones frente a S/~38,02 millones), con 11,8~\%
menos órdenes y 21,2~\% menos empresas proveedoras. Además, el gasto se concentra en pocas entidades: el
Gobierno Regional es el sector que más compra (S/~33,4 millones, 39~\% del total) y una sola entidad, el
Gobierno Regional Puno, explica el 36,1~\% del monto. Por rubros dominan las computadoras de escritorio y
portátiles, con S/~21,1 millones (24,8~\%), seguidas de los útiles de escritorio, los accesorios domésticos
y las tuberías y pinturas.

En cuanto a los proveedores, las empresas con domicilio en Puno reciben solo el 12,5~\% del monto,
mientras que las de Lima y Callao concentran el 42~\%. Asimismo, 152 de las 712 empresas, el 21,3~\%, acumulan el
80~\% del monto, un reparto muy cercano a la regla de Pareto. En el territorio ocurre algo parecido: el
distrito de Puno concentra el 71~\% del gasto (S/~60,6 millones), 6,6 veces el monto de Juliaca, que ocupa el
segundo lugar.

Los montos son muy asimétricos. La media (S/~14\,377) es 4,3 veces la mediana (S/~3\,370) y el coeficiente
de variación llega a 502~\%. Casi la mitad de las órdenes (46~\%) está entre S/~1~mil y S/~5~mil, pero el
2~\% que supera los S/~100~mil concentra el 39~\% del dinero; la orden más grande, de S/~4,13 millones, fue
una compra de computadoras de la Universidad Nacional del Altiplano. Por último, casi la totalidad de las órdenes
se emite en días hábiles: entre 1\,105 y 1\,262 por cada día de lunes a viernes, frente a 69 los sábados y
12 los domingos.

% =============================================================================
\section{Cómo se asignaron los puntajes}
\subsection{Criterios y escala}
La comparación usa doce criterios. Cinco miden la \textbf{analítica básica}, es decir, qué tan bien
permite cada herramienta preparar, calcular, describir y representar los datos. Los otros siete miden el
\textbf{\emph{reporting}}: cómo se usa, se ve, se publica y se mantiene el reporte. Se eligieron a partir
de lo que la literatura señala como esencial en un tablero, que es mostrar lo importante de un vistazo,
permitir filtrar y profundizar, y llegar a quien lo necesita \citep{yigitbasioglu2012,sarikaya2019,bach2023},
y de la manera en que el análisis visual interactivo apoya la exploración \citep{heer2012}. También se
incluyó la reproducibilidad, porque un análisis que no puede repetirse pierde valor como evidencia
\citep{peng2011,sandve2013}.

Para que el puntaje no dependiera de una apreciación general, antes de calificar se fijó qué significa cada
nivel de la escala (Tabla~\ref{tab:escala}). Cada criterio se calificó preguntando cómo resolvió la
herramienta esa tarea concreta en este proyecto, y el puntaje es el nivel cuya descripción encaja con lo
que ocurrió. Todos los criterios pesan lo mismo, así que el subtotal de cada dimensión es la suma de sus
puntajes.

\begin{table}[H]\centering\small
\caption{Escala usada para calificar cada criterio}\label{tab:escala}
\begin{tabularx}{\textwidth}{c l X}
\toprule
\textbf{Puntaje} & \textbf{Nivel} & \textbf{Qué tuvo que ocurrir en el proyecto} \\ \midrule
5 & Muy sólido & La herramienta resolvió la tarea de forma nativa o con poco esfuerzo, sin
restricciones. \\
4 & Sólido & La resolvió por completo, pero con un esfuerzo adicional moderado (configurar, aprender un
concepto nuevo, escribir algo de código). \\
3 & Aceptable & La resolvió solo en parte, mediante un rodeo, o dependiendo de una condición externa como
una licencia de pago o el permiso de un administrador. \\
2 & Débil & Solo fue posible con mucho trabajo o con conocimientos especializados previos. \\
1 & Insuficiente & No se pudo resolver en este proyecto. \\ \bottomrule
\end{tabularx}
\end{table}

\subsection{Puntajes y evidencia}
La Tabla~\ref{tab:comparacion} recoge cada puntaje junto con lo que se observó y la figura donde puede
verse. A continuación se comentan los criterios con mayor diferencia.

En \emph{estadística descriptiva} Shiny recibe 5 y Power BI 3 porque R entregó el histograma en escala
logarítmica, las densidades por sector y los percentiles sin salir de la aplicación, mientras que en Power
BI hubo que sustituir el histograma por rangos de monto y calcular cada estadístico con DAX, lo que
corresponde al nivel 3 de la escala (Figura~\ref{fig:exploracion}). En \emph{cartografía} ocurre lo
mismo: Shiny dibujó los 109 distritos con sus límites oficiales, y Power BI solo pudo representar burbujas,
además de depender de que Azure Maps esté permitido en la cuenta (Figura~\ref{fig:mapa}).

En sentido contrario, Power BI recibe 5 en \emph{interactividad} porque el filtrado cruzado, las ventanas
emergentes y el selector de métrica funcionan sin programar. Shiny recibe 4 porque todo eso es posible,
pero cada interacción, como el botón de aplicar filtros o el clic sobre el mapa, tuvo que escribirse y
probarse. La mayor brecha está en la \emph{curva de aprendizaje}: para una persona que no programa, Power BI
produce resultados desde las primeras horas de uso (5), mientras que Shiny exige saber R y entender la programación
reactiva (2). Esto coincide con la idea de la inteligencia de negocios de autoservicio, en la que el
propio usuario arma sus reportes sin depender de un programador \citep{alpar2016}, aunque incluso en ese
modelo los usuarios enfrentan vacíos de conocimiento cuando el análisis se vuelve más complejo
\citep{lennerholt2021}. En \emph{publicación} y \emph{costo} la situación se invierte. La aplicación Shiny quedó en
línea con una cuenta gratuita y un enlace abierto, mientras que compartir un reporte de Power BI con otras
personas requiere una cuenta institucional y licencias de pago, por lo que recibe 3 en ambos criterios. En
\emph{escalabilidad} Power BI obtiene 5 gracias a su motor en memoria pensado para muchos usuarios, frente a
un 3 de Shiny, cuyo plan gratuito suspende la aplicación tras un periodo de inactividad y que, con muchos usuarios a
la vez, necesitaría más servidores.

\begin{table}[p]\centering\footnotesize
\setlength{\tabcolsep}{4pt}
\renewcommand{\arraystretch}{1.25}
\caption{Puntajes por criterio y evidencia que los respalda (escala de la Tabla~\ref{tab:escala})}\label{tab:comparacion}
\rowcolors{3}{black!4}{white}
\begin{tabularx}{\textwidth}{>{\raggedright\arraybackslash}p{2.6cm} >{\raggedright\arraybackslash}X c >{\raggedright\arraybackslash}X c}
\toprule
\textbf{Criterio} & \textbf{Shiny: lo que se observó} & \textbf{Pts.} & \textbf{Power BI: lo que se observó} & \textbf{Pts.} \\ \midrule
\multicolumn{5}{l}{\cellcolor{white}\textit{Analítica básica}} \\
Preparación de datos & Transformaciones en código \texttt{dplyr}, documentadas y repetibles. & 5 &
Power Query cubre los pasos simples; la lógica queda repartida en pasos de la interfaz, más difíciles de
revisar. & 4 \\
Cálculo de indicadores & Cada indicador es una línea de R, pero hay que escribirla para cada salida. & 4 &
Medidas DAX reutilizables en todo el reporte; exigen entender el contexto de filtro. & 4 \\
Estadística descriptiva & Histograma logarítmico, densidades, percentiles y asimetría nativos
(Fig.~\ref{fig:exploracion}). & 5 & Sin histograma ni densidad; se usaron rangos de monto y estadísticos
en DAX. & 3 \\
Gráficos & \texttt{plotly} con anotaciones y control total de cada detalle. & 5 & Catálogo amplio y
pulido, con opciones limitadas por visual. & 4 \\
Cartografía & Límites oficiales de 109 distritos y detalle al hacer clic (Fig.~\ref{fig:mapa}). & 5 &
Solo burbujas por coordenadas; Azure Maps depende del administrador. & 3 \\ \midrule
\multicolumn{5}{l}{\cellcolor{white}\textit{Reporting}} \\
Interactividad & Posible, pero cada interacción se programó y probó. & 4 & Filtrado cruzado, ventana
emergente y selector de métrica sin programar. & 5 \\
Diseño & Página web con CSS libre: membrete, fondo y tarjetas propias. & 5 & Lienzo fijo y temas; buen
acabado dentro de lo que permite cada visual. & 4 \\
Curva de aprendizaje & Requiere R y programación reactiva. & 2 & Arrastrar y soltar; resultados
inmediatos. & 5 \\
Publicación & shinyapps.io gratuito y enlace público en minutos. & 5 & Compartir exige cuenta
institucional y licencia de pago. & 3 \\
Reproducibilidad & Un archivo de código versionado en GitHub. & 5 & El formato PBIP permite versionar,
aunque es menos legible que un script. & 4 \\
Costo & Software libre y plan de publicación gratuito. & 5 & Desktop gratuito; distribuir el reporte
tiene costo por usuario. & 3 \\
Escalabilidad & Con muchos usuarios se necesitan más instancias; el plan gratuito se suspende sin uso. &
3 & Motor columnar en memoria, diseñado para muchos usuarios y registros. & 5 \\ \midrule
\multicolumn{2}{l}{\cellcolor{white}\textbf{Subtotal analítica básica} (máximo 25)} & \textbf{24} & & \textbf{18} \\
\multicolumn{2}{l}{\cellcolor{white}\textbf{Subtotal \emph{reporting}} (máximo 35)} & \textbf{29} & & \textbf{29} \\
\multicolumn{2}{l}{\cellcolor{white}\textbf{Total} (máximo 60)} & \textbf{53} & & \textbf{47} \\ \bottomrule
\end{tabularx}
\end{table}

\subsection{Cómo leer el resultado}
Sumados, Shiny obtiene 53 puntos y Power BI 47, pero esa diferencia proviene íntegramente de la analítica básica. En
el \emph{reporting} empatan con 29 puntos, aunque cada una sobresale en criterios distintos: Power BI en
interactividad, facilidad de aprendizaje y escalabilidad; Shiny en diseño, publicación, reproducibilidad y
costo. Por eso el total no debe leerse como un veredicto general. Si quien mantendrá el reporte no tiene
formación en programación, el criterio de curva de aprendizaje adquiere mayor peso y la evaluación se
inclina hacia Power BI. Si lo que importa es el análisis estadístico o publicar abiertamente sin costo, favorece a
Shiny.

% =============================================================================
\section{Discusión}
Los resultados indican que ambas herramientas no compiten en el mismo ámbito. Shiny rinde mejor cuando el
análisis es el centro del trabajo: ensayar transformaciones, examinar distribuciones con escalas logarítmicas o
marcar los valores atípicos sobre el gráfico. Al ser una interfaz web sobre R, hereda todo lo que el
lenguaje ofrece, que es la razón por la que se ha vuelto común para llevar análisis especializados a
usuarios sin formación en programación \citep{jia2022}. Power BI, en cambio, está pensado para la persona
que consulta el reporte y necesita filtrar, comparar y obtener una respuesta sin tocar los cálculos, lo
que explica que se haya adoptado incluso fuera del ámbito empresarial para analizar y visualizar datos
diversos \citep{becker2019}. Las comparaciones de Power BI con otras herramientas de visualización también
suelen centrarse en la usabilidad y en las funciones de visualización para usuarios de negocio
\citep{sahaya2024}.

Esta diferencia encaja con la distinción entre explorar y comunicar. La exploración requiere libertad para
examinar los datos desde varios ángulos \citep{behrens1997}, mientras que un buen tablero debe mostrar lo
esencial de inmediato y dejar que el usuario profundice cuando lo necesite
\citep{shneiderman1996,yigitbasioglu2012}. En este proyecto el análisis exploratorio resultó más natural en
Shiny y el tablero navegable, con filtrado cruzado y ventana emergente, fue más rápido de lograr en Power
BI. La interacción, que la literatura considera central en el análisis visual \citep{heer2012}, es
justamente lo que Power BI entrega sin costo de desarrollo, mientras que en Shiny hay que construirla.

El acceso es otro factor relevante. Shiny permitió publicar el reporte de forma gratuita y abierta en poco tiempo, algo
valioso cuando el propósito es la transparencia y los datos ya son públicos \citep{janssen2012}. Power BI
está concebido para organizaciones con licencias y administradores, y publicar un reporte abiertamente
requiere habilitar opciones que muchas instituciones mantienen deshabilitadas. No es un defecto técnico, sino una
decisión de diseño orientada a controlar quién ve la información. Por eso, para una oficina de
abastecimiento cuyos usuarios necesitan consultas rápidas dentro de la entidad, Power BI sería la opción más
práctica, y para un estudio del gasto público que requiera métodos estadísticos más exigentes, Shiny ofrece
mayores ventajas.

El estudio presenta algunas limitaciones. Los puntajes reflejan la experiencia de una sola
persona con un conjunto de datos de tamaño moderado; aunque la escala se definió de antemano, sigue siendo
una valoración y no una medición con usuarios. Tampoco se probaron la actualización automática de datos, la
seguridad por filas ni un número elevado de usuarios concurrentes, escenarios en los que Power BI probablemente sumaría
ventajas. Finalmente, la región se asignó por el domicilio fiscal de cada entidad, que no siempre coincide
con el lugar donde se usa lo comprado.

% =============================================================================
\section{Conclusiones}
\begin{enumerate}
\item Con los mismos datos se construyeron dos módulos de análisis exploratorio de seis secciones cada uno,
una aplicación en Shiny y un reporte en Power BI, y ambos producen exactamente las mismas cifras, lo que
hace válida la comparación.
\item En analítica básica Shiny obtuvo más puntos (24 frente a 18), sobre todo en estadística descriptiva,
gráficos y cartografía, porque trabaja directamente sobre R.
\item En \emph{reporting} ambas herramientas obtuvieron 29 puntos. Power BI destaca por su interacción sin necesidad de programar, por
su facilidad de aprendizaje y por su escalabilidad; Shiny, por su libertad de diseño, su costo nulo y
la posibilidad de publicar abiertamente.
\item Las herramientas se complementan. Un flujo de trabajo recomendable consiste en explorar y validar en R o Shiny y trasladar
después los indicadores consolidados a Power BI cuando el reporte deba distribuirse dentro de una
organización.
\item Más allá de la comparación, el análisis aporta información relevante para la gestión regional: en 2026 el
gasto cayó 22,1~\%, se concentra en el distrito de Puno y en el Gobierno Regional, y solo el 12,5~\% del
monto se adjudica a empresas de la región.
\end{enumerate}

% =============================================================================
\newpage
\section*{Referencias}
% APA 7. Artículos de revistas indexadas en Scopus o publicaciones de IEEE, verificados en Crossref el
% 24/10/2026. Los conjuntos de datos van aparte, en Fuentes de datos.
\renewcommand{\refname}{}
\setlength{\bibsep}{6pt}
\setlength{\bibhang}{1.27cm}
\vspace{-1.2cm}
\begin{thebibliography}{99}\small
\bibitem[Alpar y Schulz(2016)]{alpar2016}
Alpar, P. y Schulz, M. (2016). Self-service business intelligence. \emph{Business \& Information Systems
Engineering, 58}(2), 151--155. \enlace{https://doi.org/10.1007/s12599-016-0424-6}{\uline{https:/}\allowbreak \uline{/}\allowbreak \uline{doi.}\allowbreak \uline{org/}\allowbreak \uline{10.}\allowbreak \uline{1007/}\allowbreak \uline{s12599-}\allowbreak \uline{016-}\allowbreak \uline{0424-}\allowbreak \uline{6}}

\bibitem[Bach et~al.(2023)]{bach2023}
Bach, B., Freeman, E., Abdul-Rahman, A., Turkay, C., Khan, S., Fan, Y. y Chen, M. (2023). Dashboard design
patterns. \emph{IEEE Transactions on Visualization and Computer Graphics, 29}(1), 342--352.
\enlace{https://doi.org/10.1109/TVCG.2022.3209448}{\uline{https:/}\allowbreak \uline{/}\allowbreak \uline{doi.}\allowbreak \uline{org/}\allowbreak \uline{10.}\allowbreak \uline{1109/}\allowbreak \uline{TVCG.}\allowbreak \uline{2022.}\allowbreak \uline{3209448}}

\bibitem[Becker y Gould(2019)]{becker2019}
Becker, L. T. y Gould, E. M. (2019). Microsoft Power BI: Extending Excel to manipulate, analyze, and
visualize diverse data. \emph{Serials Review, 45}(3), 184--188.
\enlace{https://doi.org/10.1080/00987913.2019.1644891}{\uline{https:/}\allowbreak \uline{/}\allowbreak \uline{doi.}\allowbreak \uline{org/}\allowbreak \uline{10.}\allowbreak \uline{1080/}\allowbreak \uline{00987913.}\allowbreak \uline{2019.}\allowbreak \uline{1644891}}

\bibitem[Behrens(1997)]{behrens1997}
Behrens, J. T. (1997). Principles and procedures of exploratory data analysis. \emph{Psychological
Methods, 2}(2), 131--160. \enlace{https://doi.org/10.1037/1082-989X.2.2.131}{\uline{https:/}\allowbreak \uline{/}\allowbreak \uline{doi.}\allowbreak \uline{org/}\allowbreak \uline{10.}\allowbreak \uline{1037/}\allowbreak \uline{1082-}\allowbreak \uline{989X.}\allowbreak \uline{2.}\allowbreak \uline{2.}\allowbreak \uline{131}}

\bibitem[Chen et~al.(2012)]{chen2012}
Chen, H., Chiang, R. H. L. y Storey, V. C. (2012). Business intelligence and analytics: From big data to
big impact. \emph{MIS Quarterly, 36}(4), 1165--1188. \enlace{https://doi.org/10.2307/41703503}{\uline{https:/}\allowbreak \uline{/}\allowbreak \uline{doi.}\allowbreak \uline{org/}\allowbreak \uline{10.}\allowbreak \uline{2307/}\allowbreak \uline{41703503}}

\bibitem[Fazekas y Kocsis(2020)]{fazekas2020}
Fazekas, M. y Kocsis, G. (2020). Uncovering high-level corruption: Cross-national objective corruption risk
indicators using public procurement data. \emph{British Journal of Political Science, 50}(1), 155--164.
\enlace{https://doi.org/10.1017/S0007123417000461}{\uline{https:/}\allowbreak \uline{/}\allowbreak \uline{doi.}\allowbreak \uline{org/}\allowbreak \uline{10.}\allowbreak \uline{1017/}\allowbreak \uline{S0007123417000461}}

\bibitem[Heer y Shneiderman(2012)]{heer2012}
Heer, J. y Shneiderman, B. (2012). Interactive dynamics for visual analysis. \emph{Communications of the
ACM, 55}(4), 45--54. \enlace{https://doi.org/10.1145/2133806.2133821}{\uline{https:/}\allowbreak \uline{/}\allowbreak \uline{doi.}\allowbreak \uline{org/}\allowbreak \uline{10.}\allowbreak \uline{1145/}\allowbreak \uline{2133806.}\allowbreak \uline{2133821}}

\bibitem[Janssen et~al.(2012)]{janssen2012}
Janssen, M., Charalabidis, Y. y Zuiderwijk, A. (2012). Benefits, adoption barriers and myths of open data
and open government. \emph{Information Systems Management, 29}(4), 258--268.
\enlace{https://doi.org/10.1080/10580530.2012.716740}{\uline{https:/}\allowbreak \uline{/}\allowbreak \uline{doi.}\allowbreak \uline{org/}\allowbreak \uline{10.}\allowbreak \uline{1080/}\allowbreak \uline{10580530.}\allowbreak \uline{2012.}\allowbreak \uline{716740}}

\bibitem[Jia et~al.(2022)]{jia2022}
Jia, L., Yao, W., Jiang, Y., Li, Y., Wang, Z., Li, H., Huang, F., Li, J., Chen, T. y Zhang, H. (2022).
Development of interactive biological web applications with R/Shiny. \emph{Briefings in Bioinformatics,
23}(1), Artículo bbab415. \enlace{https://doi.org/10.1093/bib/bbab415}{\uline{https:/}\allowbreak \uline{/}\allowbreak \uline{doi.}\allowbreak \uline{org/}\allowbreak \uline{10.}\allowbreak \uline{1093/}\allowbreak \uline{bib/}\allowbreak \uline{bbab415}}

\bibitem[Lennerholt et~al.(2021)]{lennerholt2021}
Lennerholt, C., van Laere, J. y Söderström, E. (2021). User-related challenges of self-service business
intelligence. \emph{Information Systems Management, 38}(4), 309--323.
\enlace{https://doi.org/10.1080/10580530.2020.1814458}{\uline{https:/}\allowbreak \uline{/}\allowbreak \uline{doi.}\allowbreak \uline{org/}\allowbreak \uline{10.}\allowbreak \uline{1080/}\allowbreak \uline{10580530.}\allowbreak \uline{2020.}\allowbreak \uline{1814458}}

\bibitem[Lu et~al.(2026)]{lu2026}
Lu, P., Jain, S., León, T. M. y White, L. A. (2026). CalCORVID: A dynamic RShiny dashboard approach to
visualize spatiotemporal clusters for public health surveillance. \emph{BMC Public Health, 26}(1),
Artículo 925. \enlace{https://doi.org/10.1186/s12889-026-26201-1}{\uline{https:/}\allowbreak \uline{/}\allowbreak \uline{doi.}\allowbreak \uline{org/}\allowbreak \uline{10.}\allowbreak \uline{1186/}\allowbreak \uline{s12889-}\allowbreak \uline{026-}\allowbreak \uline{26201-}\allowbreak \uline{1}}

\bibitem[Pebesma(2018)]{pebesma2018}
Pebesma, E. (2018). Simple features for R: Standardized support for spatial vector data. \emph{The R
Journal, 10}(1), 439--446. \enlace{https://doi.org/10.32614/RJ-2018-009}{\uline{https:/}\allowbreak \uline{/}\allowbreak \uline{doi.}\allowbreak \uline{org/}\allowbreak \uline{10.}\allowbreak \uline{32614/}\allowbreak \uline{RJ-}\allowbreak \uline{2018-}\allowbreak \uline{009}}

\bibitem[Peng(2011)]{peng2011}
Peng, R. D. (2011). Reproducible research in computational science. \emph{Science, 334}(6060),
1226--1227. \enlace{https://doi.org/10.1126/science.1213847}{\uline{https:/}\allowbreak \uline{/}\allowbreak \uline{doi.}\allowbreak \uline{org/}\allowbreak \uline{10.}\allowbreak \uline{1126/}\allowbreak \uline{science.}\allowbreak \uline{1213847}}

\bibitem[Sahaya et~al.(2024)]{sahaya2024}
Sahaya, A. P. N., Purnawarman, D., Rasyidah, F. S., Prabowo, Y. D., Sarwo y Gatc, J. (2024). Powering
sales insights: A comparative analysis of data visualization tools, Microsoft Power BI vs Tableau. En
\emph{2024 9th International Conference on Business and Industrial Research (ICBIR)} (pp. 1042--1047).
IEEE. \enlace{https://doi.org/10.1109/ICBIR61386.2024.10875740}{\uline{https:/}\allowbreak \uline{/}\allowbreak \uline{doi.}\allowbreak \uline{org/}\allowbreak \uline{10.}\allowbreak \uline{1109/}\allowbreak \uline{ICBIR61386.}\allowbreak \uline{2024.}\allowbreak \uline{10875740}}

\bibitem[Sandve et~al.(2013)]{sandve2013}
Sandve, G. K., Nekrutenko, A., Taylor, J. y Hovig, E. (2013). Ten simple rules for reproducible
computational research. \emph{PLoS Computational Biology, 9}(10), Artículo e1003285.
\enlace{https://doi.org/10.1371/journal.pcbi.1003285}{\uline{https:/}\allowbreak \uline{/}\allowbreak \uline{doi.}\allowbreak \uline{org/}\allowbreak \uline{10.}\allowbreak \uline{1371/}\allowbreak \uline{journal.}\allowbreak \uline{pcbi.}\allowbreak \uline{1003285}}

\bibitem[Sarikaya et~al.(2019)]{sarikaya2019}
Sarikaya, A., Correll, M., Bartram, L., Tory, M. y Fisher, D. (2019). What do we talk about when we talk
about dashboards? \emph{IEEE Transactions on Visualization and Computer Graphics, 25}(1), 682--692.
\enlace{https://doi.org/10.1109/TVCG.2018.2864903}{\uline{https:/}\allowbreak \uline{/}\allowbreak \uline{doi.}\allowbreak \uline{org/}\allowbreak \uline{10.}\allowbreak \uline{1109/}\allowbreak \uline{TVCG.}\allowbreak \uline{2018.}\allowbreak \uline{2864903}}

\bibitem[Shneiderman(1996)]{shneiderman1996}
Shneiderman, B. (1996). The eyes have it: A task by data type taxonomy for information visualizations. En
\emph{Proceedings 1996 IEEE Symposium on Visual Languages} (pp. 336--343). IEEE.
\enlace{https://doi.org/10.1109/VL.1996.545307}{\uline{https:/}\allowbreak \uline{/}\allowbreak \uline{doi.}\allowbreak \uline{org/}\allowbreak \uline{10.}\allowbreak \uline{1109/}\allowbreak \uline{VL.}\allowbreak \uline{1996.}\allowbreak \uline{545307}}

\bibitem[Watson y Wixom(2007)]{watson2007}
Watson, H. J. y Wixom, B. H. (2007). The current state of business intelligence. \emph{Computer, 40}(9),
96--99. \enlace{https://doi.org/10.1109/MC.2007.331}{\uline{https:/}\allowbreak \uline{/}\allowbreak \uline{doi.}\allowbreak \uline{org/}\allowbreak \uline{10.}\allowbreak \uline{1109/}\allowbreak \uline{MC.}\allowbreak \uline{2007.}\allowbreak \uline{331}}

\bibitem[Yigitbasioglu y Velcu(2012)]{yigitbasioglu2012}
Yigitbasioglu, O. M. y Velcu, O. (2012). A review of dashboards in performance management: Implications
for design and research. \emph{International Journal of Accounting Information Systems, 13}(1), 41--59.
\enlace{https://doi.org/10.1016/j.accinf.2011.08.002}{\uline{https:/}\allowbreak \uline{/}\allowbreak \uline{doi.}\allowbreak \uline{org/}\allowbreak \uline{10.}\allowbreak \uline{1016/}\allowbreak \uline{j.}\allowbreak \uline{accinf.}\allowbreak \uline{2011.}\allowbreak \uline{08.}\allowbreak \uline{002}}

\end{thebibliography}

\section*{Fuentes de datos}
\vspace{-1.2cm}
\begin{thebibliography}{99}\small
\bibitem[Central de Compras Públicas Perú Compras(2026)]{perucompras2026}
Central de Compras Públicas Perú Compras. (2026). \emph{Órdenes de compra realizadas a través de los
Catálogos Electrónicos} [Conjunto de datos]. Plataforma Nacional de Datos Abiertos.
\enlace{https://www.datosabiertos.gob.pe/dataset/órdenes-de-compra-realizadas-través-de-los-catálogos-electrónicos-central-de-compras-públicas-perú-compras}{\uline{https:/}\allowbreak \uline{/}\allowbreak \uline{www.}\allowbreak \uline{datosabiertos.}\allowbreak \uline{gob.}\allowbreak \uline{pe/}\allowbreak \uline{dataset/}\allowbreak \uline{órdenes-}\allowbreak \uline{de-}\allowbreak \uline{compra-}\allowbreak \uline{realizadas-}\allowbreak \uline{través-}\allowbreak \uline{de-}\allowbreak \uline{los-}\allowbreak \uline{catálogos-}\allowbreak \uline{electrónicos-}\allowbreak \uline{central-}\allowbreak \uline{de-}\allowbreak \uline{compras-}\allowbreak \uline{públicas-}\allowbreak \uline{perú-}\allowbreak \uline{compras}}

\bibitem[Superintendencia Nacional de Aduanas y de Administración Tributaria(2026)]{sunat2026}
Superintendencia Nacional de Aduanas y de Administración Tributaria. (2026). \emph{Padrón reducido del
RUC} [Conjunto de datos]. \enlace{https://www.sunat.gob.pe/descargaPRR/mrc137_padron_reducido.html}{\uline{https:/}\allowbreak \uline{/}\allowbreak \uline{www.}\allowbreak \uline{sunat.}\allowbreak \uline{gob.}\allowbreak \uline{pe/}\allowbreak \uline{descargaPRR/}\allowbreak \uline{mrc137\_}\allowbreak \uline{padron\_}\allowbreak \uline{reducido.}\allowbreak \uline{html}}
\end{thebibliography}

\end{document}
